%!TEX root =  CV_Master.tex
%******************************************************************
\maintitle{ACADEMIA}

\tlnew{Academic Appointments (full list)}
\cvitem{2024–}{Director, Interdisciplinary Lab (ILAS), SISSA, Trieste}
\cvitem{2023–}{\href{https://datascience.sissa.it/person/24/roberto-trotta}{Professor of Theoretical Physics} (Professore di I fascia), Physics Department, SISSA, Trieste}
\cvitem{2021–25}{Head of Data Science, SISSA, Trieste}
\cvitem{2020–22}{Associate Professor (Professore di II fascia), Physics Department, SISSA, Trieste}
\cvitem{2019–22}{\href{https://www.imperial.ac.uk/people/r.trotta}{Professor of Astrostatistics}, Imperial College London}
\cvitem{2015–20}{Director, \href{https://www.imperial.ac.uk/centre-for-languages-culture-and-communication}{Centre for Languages, Culture and Communication}, Imperial College London}
\cvitem{2016–19}{Reader in Astrophysics (equivalent to Associate Professor), Imperial College London}
\cvitem{2013–17}{Public Engagement Fellowship, Science and Technology Facilities Council, UK}
\cvitem{2012–16}{Senior Lecturer in Astrophysics, Imperial College London}
\cvitem{2008–12}{Lecturer in Astrophysics (equivalent to Assistant Professor), Imperial College London}
\cvitem{2006–08}{Junior Fellow, St Anne's College, Oxford}
\cvitem{2006}{Junior Fellow Linacre College, Oxford (declined)}
\cvitem{2005–08}{Sir Norman Lockyer Research Fellow of the Royal Astronomical Society, University of Oxford}
\cvitem{2005}{European Space Agency Research Fellowship (declined)}
\cvitem{2004–05}{Research Associate, University of Geneva, Switzerland}



\tlnew{Honours and Awards}
\cvitem{2022}{Charmian Brinson Honorary Lecture, Imperial College London}
\cvitem{2020}{Annie Maunder Medal for Outreach, Royal Astronomical Society}
\cvitem{2018}{Chaire Georges Lema\^itre, Universit\'e Catholique de Louvain, Belgium; President's Award for Leadership in Public Engagement, Imperial College London}
\cvitem{2017}{Significance Lecture, Royal Statistical Society; Outstanding Performance Award, Imperial College London; Highly commended, President's Award for Excellence in Societal Engagement, Imperial College London}
\cvitem{2016}{Longlisted, NCCPE Engage Competition Award (individual project category); President's Award for Excellence in Teaching, Imperial College London; Faculty of Natural Sciences Prize for Excellence in Teaching, Imperial College London}
\cvitem{2015}{Excellence Award, Stio town hall council, Salerno, Italy}
\cvitem{2014}{Foreign Policy 100 Global Thinkers}
\cvitem{2012}{Elected to the Royal Astronomical Society Dining Club}
\cvitem{2011}{Merit Award, Imperial College London}
\cvitem{2008}{Michelson Postdoctoral Lectureship Prize, Case Western Reserve University, USA}
\cvitem{2007}{Lord Kelvin Award Lecture, British Association for the Advancement of Science; Merit Award, Royal Astronomical Society; Merit Award, Oxford University}

\tlnew{Distinguished Fellowships and Professional Memberships}
\cvsmallheading{Elected fellow.}
European Coalition for AI in Fundamental Physics (EuCAIF, 2025); International Astrostatistics Association (2016); Royal Astronomical Society (FRAS); Royal Astronomical Society Dining Club (2012).

\cvsmallheading{Academic / senior fellow.}
Imperial College London Data Science Institute (2018–22); Higher Education Academy, UK (Senior Fellow, 2021; Fellow status achieved 2011).

\cvsmallheading{Professional member.}
International Society for Bayesian Analysis; International Astronomical Union; Swiss Society for Astronomy and Astrophysics; European Astronomical Society; Gresham Society.

\cvsmallheading{Affiliate.}
Chartered Management Institute, UK (2021–2022).

\tlnew{Continuing Professional Development}
\cvitem{2015–21}{Imperial Leadership and Management Development Programme (validated by the Chartered Management Institute, UK). Courses on: introduction to management; managing groups and teams; leading successful teams; effective recruitment and selection; finance for non-finance managers; promoting equality and diversity; CPD portfolio; 360 review and action planning; decision making; unconscious bias; managing difficult situations.}
\cvitem{2016}{Introduction to IP rights}
\cvitem{2014}{Bond Solon Expert Witness training programme}

\tlnew{Scientific Editor}
\cvitem{2021–}{Scientific Editor and member of the inaugural Editorial Board of \href{https://academic.oup.com/rasti}{\em RAS Techniques and Instruments} (RASTI), the open-access journal of the Royal Astronomical Society covering data collection methods, instrumentation, data processing and modelling, AI and machine learning, statistical methods, software, and high-performance computing.}

\tlnewnospace{Scientific Activities}

\cvitem{2026–}{Management Committee Member, EU Cost Action `Artificial Intelligence for Fundamental Physics: Open Roadmapping and Inclusive Actions (AIFORIA)'}
\cvitem{2025–}{Member, LSST Dark Energy Science Collaboration}
\cvitem{2022–}{Scientific Consultant, International Centre for Theoretical Physics (ICTP), Trieste; Programme Committee Member, `ML and AI for Science', FVG Data Science Institute; Member, XLZD Consortium for direct dark matter detection}
\cvitem{2014–}{DARWIN direct dark matter detection experiment: Executive Board member (2015–), Science Working Package leader (2014–)}
\cvitem{2014–16}{ATLAS Collaboration, Short Term Associate}
\cvitem{2014–18}{XENON100 direct dark matter detection experiment, Associated Scientist}
\cvitem{2012–15}{Fermi-LAT Affiliated Scientist}
\cvitem{2012, 2013}{Guest editor: \emph{Statistical Analysis and Data Mining} special issue (2012); \emph{Physics of the Dark Universe} special issue (2013)}
\cvitem{2010–14}{Euclid ESA space mission, Theory Working Group member (2010–14); likelihood analysis Work Package co-lead (2020–22)}
\cvitem{2005–07}{Scientific Advisor on Dark Energy probes, STFC, UK}
\cvitem{2006–15}{Co-author of the \texttt{SuperBayeS} public code for the analysis of supersymmetric theories}
\cvitem{2001–}{30+ invited scientific talks, 110+ research talks and seminars}

\tlnewnospace{Service, Advisory and Evaluation}
\cvitem{2026–27}{External Advisory Board for the McWilliams Center for Cosmology and Astrophysics at Carnegie Mellon University, USA} 
\cvitem{2026}{Habilitation panel, U. Nova Gorica, Slovenia. }
\cvitem{2026}{Agence Nationale de la Recherche (France) evaluator; Dutch Research Council evaluator; ERC Consolidator Grants evaluator}
\cvitem{2025}{Emmanuel College Cambridge Research Fellowship assessor; Cambridge University promotion assessor; UKRI Future Leaders evaluator; ERC Synergy Grants evaluator; Swiss National Science Foundation evaluator}
\cvitem{2024}{Dutch Research Council evaluator; Swiss National Science Foundation evaluator}
\cvitem{2023}{Tenure committee reviewer, School of Statistics, University of Minnesota}
\cvitem{2021}{ERC Starting Grants evaluator}
\cvitem{2020}{Referee, Canada Research Chairs programme}
\cvitem{2018}{Member, STFC Public Engagement Early Career Researcher Network Review Panel}
\cvitem{2016–}{Member, Working Group on Equity and Inclusion, Commission C1 of the IAU}
\cvitem{2015–}{Member, DARWIN-LXe Executive Board; Outreach and Knowledge Transfer Leader, International Astrostatistics Association; Member, Council of the International Astrostatistics Association}
\cvitem{2015}{University Research Fellowships referee, Royal Society; Imperial College London representative, Interdisciplinary Curriculum Group}
\cvitem{2014–15}{STFC Rutherford Fellowships panel}
\cvitem{2011–}{Member, Astrostatistics Committee, International Statistical Institute}
\cvitem{2010}{Founding member, Imperial Centre for Inference and Cosmology}
\cvitem{2009}{Jesus College and Peterhouse College fellowships referee, Cambridge}
\cvitem{2009–2012}{Member, \emph{Astronomy \& Geophysics} Editorial Board}
\cvitem{2007–}{South African National Research Foundation evaluator}
\cvitem{2007–12}{Royal Society International Grants Panel}
\cvitem{Ongoing}{Referee for \emph{Nature Physics}, \emph{PRL}, \emph{ApJ}, \emph{JCAP}, \emph{JHEP}, \emph{PLB}, \emph{Phys.~Rev.~D/E}, \emph{New Astron.}, \emph{MNRAS}, \emph{Comp.~Phys.~Comm.}, \emph{Science and Education}, \emph{RASTI} --- 80+ papers refereed.}


\tlnewnospace{Institutional Responsibilities, Management and Leadership}

\cvsubsection{SISSA, Trieste}
\cvitem{2024–}{Director, Interdisciplinary Lab (ILAS); member, Library committee}
\cvitem{2021–25}{PhD Coordinator and Academic Board Chair, Theoretical and Scientific Data Science; Head of Group, Theoretical and Scientific Data Science; Member, Giunta dell'Area Fisica (Physics Area Steering Board)}
\cvitem{2021–}{Member, Laboratorio Interdisciplinare per le Scienze Naturali (ILAS); Physics Area representative, Knowledge Transfer Committee (Commissione per la Valorizzazione del Talento e il Trasferimento della Conoscenza)}
\cvitem{2021}{Physics Area coordinator for the 2014–19 national assessment exercise (VQR)}

\cvsubsection{Gresham College, London}
\cvitem{2019–22}{Member of the Academic Board}

\cvsubsection{International}
\cvitem{2024–}{Member, EuCAIF Steering Board, European Coalition for AI in Fundamental Physics}

\cvsubsection{Imperial College London}
\cvitem{2020}{Education Strategy and Operations Group}
\cvitem{2019–20}{Horizons Stream Lead, I-Explore programme}
\cvitem{2020}{Education Office COVID-19 Response Group}
\cvitem{2018–19}{Chair, Curriculum Review Reference Panel}
\cvitem{2015–20}{Chair, CLCC Departmental Teaching Committee}
\cvitem{2016–20}{Chair, President's Award review panel `Supporting the Student Experience'}
\cvitem{2015–20}{Chair, Imperial Horizons Team; Chair, CLCC Management Team}
\cvitem{2016–20}{Co-chair, CLCC/CHERS Education Committee}
\cvitem{2018–20}{Member, I-Explore Modules Innovation Group; Member, Enterprise Academic Advisory Group; Member, Engagement Strategy Group}
\cvitem{2015–20}{Member, Programmes Committee}
\cvitem{2013–20}{Member, Imperial SpaceLab Steering Group}
\cvitem{2016–20}{Member, Strategy Board, Centre for Performance Science}
\cvitem{2017–20}{Member, Online Learning Innovation Group}
\cvitem{2016–18}{Member, Academic Partnerships Engagement with Research Steering Group}
\cvitem{2015–20}{Member, Education Office Management Committee}
\cvitem{2016–20}{Member, Imperial College Senate}
\cvitem{2015–18}{Member, Imperial Horizons Committee}
\cvitem{2015–17}{Member, Task and Finish Group for Assessment}

\tlnewnospace{Organising of Scientific Meetings}
\setcvcounter{nOrgMeetings}
\cvcount{2027}{SOC and co-proposer, `AI in Astronomy Education and Outreach: Opportunities and Risks', IAU General Assembly, Rome, Italy}
\cvcount{2026}{SOC, `AI in Education', SISSA, Trieste, Italy; SOC, `DESC Virtual Collaboration Meeting' (online)\cvsep LOC, `Quasar Bazar', IFPU, Trieste\cvsep SOC, Advanced School in Applied Machine Learning: Uncertainty and Robustness of DL Models for Science, ICTP \cvsep SOC and co-proposer, NeurIPS workshop `AI and Science: Evolution or Extinction?'\cvsep SOC, INSAP XIV conference, Rome}
\cvcount{2025}{LOC Chair, Convegno Nazionale di Comunicazione della Scienza, Trieste\cvsep SOC, Advanced School on Foundation Models for Scientific Discovery, ICTP, Trieste\cvsep SOC, AI \& Physics track, `AI4Science' conference, Ljubljana\cvsep SOC, AI Alliance Meet-up, Trieste\cvsep SOC, INSAP XIII conference, Belfast\cvsep SOC, EuCAIFCon 2025, Cagliari}
\cvcount{2024}{SOC, Advanced School on Applied Machine Learning, ICTP, Trieste\cvsep SOC, StatPhys conference, Imperial College London\cvsep SOC, EuCAIFCon 2024, Amsterdam\cvsep SOC, INSAP XII conference, Corfu\cvsep LOC, `Unsupervised Exploration of Fast Radio Bursts Dynamic Spectra', IFPU Team Research Workshop, Trieste\cvsep SOC, `Statistics meets ML', Imperial College London\cvsep SOC, COSMO21 conference, Chania}
\cvcount{2023}{Organiser, `AI Interactions: ChatGPT and Beyond', Trieste and online\cvsep SOC, `Machine Learning Enabled Searches for New Physics in Astrophysical Data', IFPU, Trieste}
\cvcount{2022}{Convenor and moderator, Year of Basic Sciences for Sustainable Development: Embedding Ethics in Machine Learning, ICTP, Trieste (online)\cvsep SOC, `Ethical and Societal Challenges of Machine Learning', ICTP, Trieste (online)\cvsep SOC, `Mind the Gap: How We Build, Perceive, Experience and Share Spaces', SISSA\cvsep SOC, INSAP XI, CalTech}
\cvcount{2017–}{International Executive Committee, `Inspiration of Astronomical Phenomena' (INSAP)}
\cvcount{2018}{SOC, Astronomical Data Analysis 9, Valencia}
\cvcount{2016}{SOC, Astronomical Data Analysis 8, Crete}
\cvcount{2015}{Co-Chair, `Identification of Dark Matter with a Cross-Disciplinary Approach' workshop, UAM Madrid}
\cvcount{2014}{Chair, `From Cosmology to Customers: Astrostatistics Solutions to Big Data Problems' industry–academia workshop, Imperial}
\cvcount{2009, 2013}{Chair and founder, `Cosmostat' conference series: Centro Stefano Franscini, Switzerland (2009)\cvsep Banff International Research Station (2013)}
\cvcount{2013}{Lead Organiser, `Multi-probe Identification of Dark Matter', Kavli ITP programme, UCSB}
\cvcount{2012}{Chair, DarkAttack12, Centro Stefano Franscini, Switzerland\cvsep SOC, ICIC inaugural workshop, Imperial College London\cvsep SOC, `Statistical Issues in Searches', SLAC, Palo Alto}
\cvcount{2010}{SOC, RAS National Astronomy Meeting, Glasgow\cvsep Co-chair, RAS Specialist Discussion Meeting `Novel Methods for the Exploitation of Large Astronomical and Cosmological Data Sets', London\cvsep Co-chair, IOP dark matter meeting, London\cvsep Co-chair, Astrostatistics workshop, London}
\cvcount{2012, 2013}{Founder and chair, `The Big Questions' public lecture series (2012) and `The Sensual Universe' public lecture series (2013), London}
\cvcount{2008}{Co-chair, `Bayesian Evidence' workshop, Brighton}
\cvcount{2007}{Co-chair, RAS Specialist Discussion Meeting `Statistical Challenges in Particle Astrophysics and Cosmology', London}
\cvcount{2002–05}{Co-chair, Journ\'ees de Cosmologie des Lacs Alpins, Geneva/Annecy/Lausanne/CERN}

\labelcvcount{nOrgMeetings}
\clearcvcounter
\tlnewnospace{Teaching}

\cvsubsection{Undergraduate courses and lectures}
\cvitem{2021}{Guest lunchtime lecture, `The Machine Learning Revolution in Cosmology', Imperial (online)}
\cvitem{2015–20}{Lecturer, `Physics of the Universe', 3rd-year physics core course, Imperial}
\cvitem{2008–12}{Lecturer, `Statistics of Measurement', 2nd-year physics core course, Imperial}
\cvitem{2010–20}{Associate lecturer: `Cosmology', Imperial}
\cvitem{2008–09}{Associate lecturer: `Sun, Stars and Planets' (2008–09)}
\cvitem{2017, 2018}{Guest lecturer, `Communicating Science', Imperial Horizons course}
\cvitem{2016}{Guest lecturer, `Creative Writing', Imperial Horizons course}
\cvitem{2010}{Guest lecturer, `Astrophysics', 4th-year option course, Imperial}

\cvsubsection{Postgraduate courses at SISSA}
\cvitem{2023–}{Bayesian Methods I, Theoretical and Scientific Data Science PhD programme}
\cvitem{2020–}{Ethical Aspects of AI, Theoretical and Scientific Data Science PhD programme}
\cvitem{2021–}{Intelligenza Artificiale e Informazione (with Elisabetta Tola), Master in Science Communication}
\cvitem{2020–23}{Bayesian Methods II, Theoretical and Scientific Data Science PhD programme}
\cvitem{2021–}{Bayesian Inference and Machine Learning in Cosmology, Theoretical and Scientific Data Science PhD programme}

\cvsubsection{Postgraduate lectures, masterclasses and graduate schools}
\cvitem{2026}{AI training course, SISSA; EFC network AI training workshop, Bertinoro}
\cvitem{2024–}{Machine Learning for Business Analytics, MBA in International Business, MIB Trieste}
\cvitem{2023}{Introduction to Bayesian Inference, SISSA Master in HPC, Trieste; Machine Learning for Business Analytics, MBA in International Business, MIB Trieste}
\cvitem{2022}{Introduction to Data Science, International Master in Physics of Complex Systems, ICTP, Trieste; Statistics lectures, Cosmology Summer School, ICTP, Trieste}
\cvitem{2021}{Avventure fra Arte e Scienza nel Comunicare la Cosmologia, SISSA MSC lecture; From Bayes to Machine Learning, Heidelberg Grad Days (online); Bayesian Inference, Saas-Fee Course on `Astronomy in the Era of Big Data' (online)}
\cvitem{2019}{Keynote, Research Computing Summer School, Imperial College London}
\cvitem{2018}{ICIC Data Analysis Postgraduate School, Imperial College London; Analytics, Inference and Computation in Cosmology: Advanced Methods, Carg\`ese, France; Astronomical Data Analysis IX Summer School, Valencia, Spain}
\cvitem{2017}{Stats4Astro Astrostatistics School, Autrans, France; ISAPP Summer School, Texel, The Netherlands; First Italian Astrostatistics School, Milan}
\cvitem{2016}{European Space Astronomy Centre (ESAC) Statistics Workshop, Villanueva de la Ca\~nada, Spain; ADA8 Summer School on Astronomical Data Analysis, Chania, Greece}
\cvitem{2014}{XII School of Cosmology, Carg\`ese, France; Niels Bohr Institute PhD School, `Statistical Methods', Copenhagen; XXVI Winter School, `Bayesian Inference in Astronomy and Astrophysics', Canary Islands; ICIC STFC Data Analysis School, London; 44th Saas-Fee Advanced Course on Astronomy and Astrophysics, `Cosmology with Wide-Field Surveys', Switzerland}
\cvitem{2013}{School on Bayesian Analysis in Physics and Astronomy, Stellenbosch, South Africa; ICIC Data Analysis Postgraduate School, London}
\cvitem{2012}{`Statistical Inference', African Institute for Mathematical Sciences postgraduate course, Cape Town}
\cvitem{2011}{`Advanced Statistical Methods for Cosmology and Astroparticle Physics', Niels Bohr Institute, Copenhagen}
\cvitem{2010}{Cosmology postgraduate lectures, Imperial; `Advanced Statistical Methods' lectures, Imperial; `Statistical Methods for Cosmology', 1st Jayme Tiomno School of Cosmology, Rio de Janeiro}
\cvitem{2009}{`Advanced Statistical Tools for Cosmology', Valencia University; `Probability Theory, Statistics and Data Analysis' postgraduate lectures, Imperial/UCL; `The CMB and its Statistical Interpretation', W\"urzburg University postgraduate lecture series; `Modern Cosmology' postgraduate lectures, Imperial/UCL; `Advanced Statistical Tools for Cosmology' postgraduate lectures, U. Aut\'onoma Madrid}
\tlnewnospace{Supervision}

\cvsubsection{Postdoctoral researchers}
\cvitem{2025–}{Giulio Scelfo}
\cvitem{2024–26}{Junsong Cang}
\cvitem{2024–25}{Martin de los Rios (now faculty at CONICET, Argentina)}
\cvitem{2024–25}{David Prelogovi\'c}
\cvitem{2023–}{Rahul Srinivasan}
\cvitem{2023–25}{Chiara Moretti}
\cvitem{2023}{Matteo Breschi}
\cvitem{2022–}{Andre Scaffidi}
\cvitem{2022–24}{Gabriella Contardo}
\cvitem{2020–22}{Ira Wolfson}
\cvitem{2020–21}{Victor Bonjean}
\cvitem{2018–20}{Eliel Camargo-Molina (now faculty at U. Uppsala)}
\cvitem{2017–20}{Alex Geringer-Sameth (now permanent staff at Lawrence Livermore National Lab)}
\cvitem{2014}{Kaisey Mandel (now faculty at Cambridge)}

\cvsubsection{Current PhD and MSc students}
\cvitem{2023–}{M. Giunta, G. Puleo}

\cvsubsection{Completed PhD students (thesis title; distinctions and first destination)}
\cvitem{2022–26}{M. Rigo (SISSA), \emph{Towards fast and robust Large Scale Structure analysis with physics-based emulators}}
\cvitem{2022–26}{S. Mishra (SISSA), \emph{Generating the Universe: Field-Level Diffusion Models for Emulation, Reconstruction, and Posterior Inference} – awarded {\em cum laude}} 
\cvitem{2021–25}{M. Geng (SISSA), \emph{Reflections on Distributions: Humans, Machines, Languages} (postdoc at \'Ecole Normale Sup\'erieure, Paris)}
\cvitem{2020–24}{K. Karchev (SISSA), \emph{SNIa Cosmology for the 21st Century} – awarded {\em cum laude}; IAU  Best Thesis Prize: Facilities, Technologies and Data Science 2025; Springer Thesis Prize 2025 (Simons Foundation Fellow, Barcelona)}
\cvitem{2020–23}{M. Autenrieth (Imperial Statistics, co-supervision), \emph{Principled Bayesian Modelling and Statistical Learning with Non-Representative Data in Astrophysics} – Royal Statistical Society Emerging Applications Thesis Prize 2025 (postdoc at Imperial)}
\cvitem{2017–22}{W. Rahman (Imperial), \emph{Constraining the Anisotropic Expansion of the Universe with Type Ia Supernovae and Improving the Treatment of Selection Effects within Bayesian Hierarchical Models} (data scientist at Meta)}
\cvitem{2015–19}{S. Hoof (Imperial, co-supervision), \emph{Global Fits of Axions and WIMPs in Astrophysics, Cosmology and Particle Physics} (Imperial President Scholarship holder; postdoc at G\"ottingen)}
\cvitem{2014–18}{J. McKay (Imperial, co-supervision), \emph{Constraining Dark Matter with Renormalisation and Global Fits} (Imperial President Scholarship holder; software developer at Verizon Connect)}
\cvitem{2013–17}{H. Shariff (Imperial), \emph{Application of Bayesian Statistics in Supernovae Ia Cosmology} (investment consultant at Mercer)}
\cvitem{2011–14}{C. Strege (Imperial), \emph{Characterisation of Particle Dark Matter via Multiple Probes} --- Winton Capital best PhD thesis in physics 2014 (quant at CityBank)}
\cvitem{2008–11}{M. March (Imperial), \emph{Advanced Statistical Methods for Astrophysical Probes of Cosmology} --- Springer Thesis Prize 2011, Winton Capital best PhD thesis in physics 2011 (postdoc at Penn State)}

\cvsubsection{MSc / MSci students}
\cvitem{2025}{U. Zivanovi\'c (U. Trieste), \emph{An Application of the Rotary Masked Autoencoder to Irregular Sparse Astrophysics Datasets}}
\cvitem{2023}{R. Serra (U. Trieste), \emph{Improving Photometric Galaxy Redshift Estimation under Covariate Shift with StratLearn}}
\cvitem{2022}{S. di Gioia (U. Trieste), \emph{Fast and Parallel Algorithms for Causal Discovery in Astrophysics}}
\cvitem{2021}{R. Corti (U. Trieste), \emph{Machine Learning Methods for Supernova Type Ia Selection Effects Modelling}}
\cvitem{2019}{M. von Wietersheim-Kramsta (Imperial), \emph{Measuring the Galaxy Mass Function Using Photometric Data}}
\cvitem{2017}{H. Bouvier (Imperial), \emph{Determining Type Ia Supernovae Host Galaxy Characteristics Using Galaxy Zoo: Hubble and the \textsc{Prospector}-$\alpha$ Model}}
\cvitem{2017}{S. Kobayashi (Imperial), \emph{Determining Type Ia Supernovae Host Galaxy Characteristics Using Galaxy Zoo: Hubble and the \textsc{Prospector}-$\alpha$ Model}}
\cvitem{2016}{E. Revsbech (Imperial Statistics, co-supervision), \emph{Synthetically Augmented Light Curve Classification} --- Winton Capital Prize Best Project; I. Siska (Imperial), \emph{Constraining Anisotropic Cosmological Expansion using Supernovae Type Ia within a Bayesian Statistical Framework}; A. Monge Imedio (Imperial), \emph{Constraining Cosmological Anisotropies with SNIa Data Using Bayesian Methods}}
\cvitem{2015}{S. Pak (Imperial), \emph{Constraining Cosmic Expansion Anisotropy using Type Ia Supernovae}; K. Blanchette (Imperial), \emph{Cosmological Parameter Inference in a Bayesian Hierarchical Model with Redshift-Dependent Hubble Residuals}}
\cvitem{2011}{G. Pisani (co-supervision, Rome La Sapienza)}

\cvsubsection{Other supervision}
\cvitem{2011–20}{35 BSc projects, Imperial}
\cvitem{}{22 Summer interns / UROP / IROP / ICTP / UWC undergraduate students supervised}

\cvsubsection{Mentorship}
\cvitem{2026–}{Active member, LSST/DESC mentoring programme}

\tlnewnospace{Examining at Postgraduate Level}



\cvsubsection{External examiner / assessor}
\cvitem{2026}{B. Ricketts (PhD, U. of Amsterdam)}
\cvitem{2026}{A. Bairagi (PhD, Institut d'Astrophysique de Paris)}
\cvitem{2025}{V. Blasone (PhD, U. Trieste, Applied Data Science and AI programme)}
\cvitem{2024}{Z. Rokavec (MSc, U. Nova Gorica); A. Prajapati (PhD, Gran Sasso Science Institute); D. Prelogovi\'c (PhD, Scuola Normale Superiore; committee member); S. Rinaldi (PhD, U. Pisa)}
\cvitem{2023}{D. Lanzieri (PhD, CEA Paris-Saclay)}
\cvitem{2022}{J. Thorp (PhD, Cambridge); F. Agocs (PhD, Cambridge)}
\cvitem{2019}{R. Hardwick (PhD, U. Portsmouth)}
\cvitem{2018}{J. Faulkner (PhD, U. Edinburgh)}
\cvitem{2017}{S. Hee (PhD, U. Cambridge); H. Mootoovaloo (MSc, U. Cape Town)}
\cvitem{2016}{T. Walwyn (MSc, U. Cape Town)}
\cvitem{2014}{F. Vansyngel (PhD, Institut d'Astrophysique de Paris)}
\cvitem{2013}{H. Silverwood (MSc, U. Canterbury, NZ)}
\cvitem{2011}{Y. Akrami (PhD, Stockholm University; opponent); J. Newling (MSc, U. Cape Town)}

\cvsubsection{Internal examiner, Imperial College London}
\cvitem{2022}{E. Hauke (PhD, Centre for Higher Education Research and Scholarship)}
\cvitem{2017}{Shijing Si (PhD, Statistics Section)}
\cvitem{2016}{J. Alsing (PhD, Astrophysics group)}
\cvitem{2013}{H. Shariff (MSc, Astrophysics group)}
\tlnewnospace{Invited Scientific Talks}
\setcvcounter{nInvitedTalks}
\cvcount{2026}{Plenary talk, EUCAIFCon, Heidelberg \cvsep Physics dept, Ljubljana U. \cvsep Data science seminar, CERN}
\cvcount{2025}{Panelist, `Agentic AI, LLMs and Research', SISSA, Trieste\cvsep Keynote, IFPU Focus Week, Trieste\cvsep Keynote, Randstad Foundation think tank meeting, Milan\cvsep Keynote, Ital-IA, Italian National meeting on AI, Trieste}
\cvcount{2024}{Panel, Climate Conference Arena, Bologna\cvsep Colloquium, Celebrating 60 Years of the Associates Programme, ICTP, Trieste\cvsep AI panel moderator, ICTP Symposium `The Future of Scientific Computing', Trieste\cvsep \href{https://www.youtube.com/watch?v=PPWYvDpY8Bk&ab_channel=UniverzavNoviGorici\%2FUniversityofNovaGorica}{Scientific Afternoons of the University of Nova Gorica}, Nova Gorica\cvsep CMSP News and Views Colloquium, ICTP, Trieste\cvsep Erwin Schr\"odinger Colloquium, University of Z\"urich\cvsep Keynote, Research in Dialogue, University of Bern\cvsep Heidelberg Joint Astronomical Colloquium}
\cvcount{2023}{Astrophysics Colloquium, Ludwig-Maximilians-Universit\"at, Munich\cvsep Colloquium Series in Theoretical and Computational Physics (CSTCP)}
\cvcount{2022}{`StratLearn: A General-Purpose Statistical Method for Improved Learning under Covariate Shift', IAA/IAU Astroinformatics and Astrostatistics Seminars (online)}
\cvcount{2021}{`The Impact of Machine Learning in Cosmology' (online)}
\cvcount{2019}{Plenary, Picturing the Invisible Conference, UAL, London\cvsep Keynote, `Artificial Intelligence: Art or Science?', SISSA, Trieste\cvsep Keynote, ICTP, Trieste\cvsep AHRC Network meeting}
\cvcount{2017}{Colloquium, University of Bern\cvsep Colloquium, Indian Institute of Technology Madras}
\cvcount{2016}{Keynote, Global Challenges Showcase event, Imperial College London\cvsep Keynote, COSMO21 Conference, Chania, Greece}
\cvcount{2015}{Southampton Physics Colloquium\cvsep BASP Frontiers 2015, Villars, Switzerland}
\cvcount{2013}{Director's Blackboard Talk, Kavli Institute, Santa Barbara}
\cvcount{2012}{Invited plenary, Statistical Issues in Searches, SLAC, Stanford\cvsep Invited plenary, Astronomical Data Analysis 7, Carg\`ese\cvsep Invited plenary, MaxEnt 2012, Garching}
\cvcount{2011}{Invited plenary, 58th ISI World Congress, Dublin\cvsep Invited plenary, Statistical Challenges in Modern Astronomy V, Penn State\cvsep Invited talk, IOP dark matter meeting, London\cvsep Invited plenary, Aspen Winter Conference, `Indirect and Direct Detection of Dark Matter', Aspen\cvsep Colloquium, AIMS, Cape Town}
\cvcount{2010}{Colloquium, EPFL\cvsep Invited plenary, PROSPECT workshop, Stockholm\cvsep Summary talk, `Statistical Issues Relevant to Significance of Discovery Claims', Banff\cvsep Invited plenary, TeV Particle Astrophysics, Paris}
\cvcount{2009}{Review talk, ILIAS/ENTApP meeting, CERN, Geneva\cvsep Invited talk, National Astronomy Meeting, Hatfield}
\cvcount{2008}{Michelson Prize Lecture, Case Western Reserve University, USA\cvsep Invited plenary, ESF Exploratory Workshop, Oporto\cvsep Invited talk, Tools08, Munich\cvsep Review talk, XXVI Workshop on Recent Developments in High Energy Physics and Cosmology, Ancient Olympia}
\cvcount{2007}{Lord Kelvin Award Lecture\cvsep Invited talk, TeV Particle Astrophysics 2007, Venice}
\cvcount{2006}{Invited talk, 11th Marcel Grossmann Meeting, Berlin\cvsep Invited plenary, RAS General Meeting, London}
\cvcount{2005}{Review talk, `Recent Advances in Particle Physics and Cosmology', Thessaloniki}

\vspace{\baselineskip}
\labelcvcount{nInvitedTalks}
\clearcvcounter
\tlnewnospace{Contributed Conference Talks}
\setcvcounter{nContribTalks}
\cvcount{2024}{The Dark Matters, parallel session at the RGS-IBG Annual International Conference (online)}
\cvcount{2021}{Debating the potential of machine learning in astronomical surveys (online)}
\cvcount{2019}{CosmoAI, Ascona, Switzerland}
\cvcount{2018}{Cosmo21, Valencia}
\cvcount{2017}{The Road to the Stars, Santiago de Compostela}
\cvcount{2016}{XENON100 Collaboration Meeting, Mainz, Germany}
\cvcount{2013}{ISI World Statistics Meeting, Hong Kong}
\cvcount{2012}{ISBA World Meeting, Kyoto}
\cvcount{2011}{ERC Kickoff meeting, Amsterdam\cvsep IOP Dark Matter Meeting at King's College\cvsep 58th ISI World Congress\cvsep PHYSTAT2011, CERN, Geneva}
\cvcount{2009}{Corfu Summer Institute, Greece\cvsep cosmostats09, Ascona, Switzerland\cvsep Galileo Galilei Institute dark energy conference, Florence}
\cvcount{2008}{Nested Sampling workshop, Cambridge\cvsep Bayesian evidence workshop, University of Sussex\cvsep CosmoTools08, Marseille}
\cvcount{2007}{Initial Conditions in Cosmology workshop, W\"urzburg, Germany\cvsep CERN dark matter workshop, Geneva}
\cvcount{2006}{RAS Young Astronomers Meeting, London\cvsep Identification of Dark Matter Conference, Rhodes, Greece\cvsep Galileo Galilei Institute Cosmology Workshop, Florence\cvsep The Dark Side of the Universe Workshop, Madrid\cvsep CosmoUK Meeting, Oxford\cvsep Francesco Melchiorri Memorial Conference, Rome\cvsep TeV Particle Astrophysics II, Madison, Wisconsin\cvsep Astroparticle UK Meeting, Sheffield\cvsep National Astronomy Meeting, Leicester\cvsep SKA conference, Oxford\cvsep St Anne's College, Oxford}
\cvcount{2005}{European Dark Energy Network Meeting, Paris\cvsep Subaru/Gemini dark energy meeting, Waikoloa, Hawaii\cvsep PHYSTAT05, Oxford\cvsep Sils Maria Cosmology Workshop, Engadin\cvsep ENTApP meeting, CERN, Geneva}
\cvcount{2004}{7th Journ\'ee des Lacs Alpins de Cosmologie, EPFL Lausanne\cvsep 6th Journ\'ee de Cosmologie des Lacs Alpins, Geneva}
\cvcount{2003}{Second Aegean Summer School on the Early Universe, Syros\cvsep Sils Maria Cosmology Workshop, Engadin\cvsep ATLAS Collaboration Physics Tutorial, CERN\cvsep X Marcel Grossmann Meeting, Rio de Janeiro\cvsep CMBNet Workshop on Science and Parameter Extraction, Oxford}
\cvcount{2002}{5th Journ\'ee de Cosmologie des Lacs Alpins, Annecy\cvsep CMBNet Workshop, Geneva\cvsep Santa Fe Cosmology Workshop, New Mexico}
\cvcount{2001}{CMBNet General Meeting, Frascati, Italy\cvsep UK Cosmology Meeting, Ambleside}

%
%\vspace{\baselineskip}
%\tlnewnospace{Research Seminars} 
%\begin{enumerate}
%\item The future of Education, CLCC Research Seminar, Nov 2019 
%\item March 19: Talk to DSI Advisory Board 
% \item March 19: DSI Fellows Event Pecha Kucha 
%\item Jan 19: SISSA Data Science talk
%\item Nov-16 Centre for Languages, Culture and Communication Research Seminar (2016)
%\item KITP Director's Blackboard talk, Santa Barbara, UCSB (2013)
%\item Sep-12	CCAPP, Ohio State University
%\item Feb-12	GRAPPA Amastedam
%\item Jan-12	Bonn U.
%\item Sep-11	Univeristy of Zurich
%\item Apr-11	AIMS
%\item Dec-10	Queen's  University, Belfast
%\item Nov-10	EPFL 
%\item Nov-09	University of Zurich
%\item Nov-09	King's College London
%\item Oct-09	Royal Holloway, London
%\item Sep-09	Uppsala University
%\item Sep-09	Centre for Dynamical Processes and Structure Formation, Uppsala University, Sweden
%\item Jun-09	Osservatorio Astronomico di Trieste, 
%\item Jun-09	Mullard Space Science Laboratory, Holmbury St. Mary
%\item May-09	IFIC, Valencia
%\item May-09	Osservatorio Astronomico di Brera, Milan,
%\item Jan-09	Royal Observatory Edinburgh
%\item May-08	UC Irvine
%\item May-08	Liverpool John Moores University
%\item May-08	JPL, Pasadena
%\item May-08	CWRU 
%\item May-08	CWRU
%\item May-08	CWRU
%\item Jan-08	Oxford Philosophy of Physics Dept
%\item Jan-08	Theoretical Physics Dept, Universite Catholique de Louvain, Belgium
%\item Jan-08	Oxford University
%\item Nov-07	Theoretical Physics Dept, Heidelberg, Germany
%\item Oct-07	University of Geneva
%\item Oct-07	Imperial College London
%\item Jul-07	Universita di Bari
%\item May-07	Royal Statistical Society, Oxford
%\item May-07	Cavendish Lab, Cambridge
%\item Apr-07	Dark cosmology center, Copenhagen, Denmark
%\item Apr-07	Niels Bohr Institute, Copenhagen, Denmark
%\item Mar-07	Fermilab
%\item Mar-07	Perimeter Institute for Theoretical Physics, Waterloo, Canada
%\item Mar-07	CWRU
%\item Mar-07	Kavli Institute for Cosmological Physics, Chicago
%\item Mar-07	Universita di Salerno
%\item Feb-07	Queen Mary, London
%\item Jun-06	University of Sussex
%\item May-06	UCL
%\item May-06	Oxford Astrophysics 
%\item Apr-06	DAMPT, Cambridge
%\item Jan-06	Oxford Astrophysics
%\item Jan-06	Laboratoire de Physique Theorique, Orsay
%\item Jan-06	RWTH Aachen, Germany
%\item Dec-05	Sheffield University
%\item Dec-05	Royal Observatory Edinburgh
%\item Oct-05	Lancaster University
%\item Oct-05	Oxford Astrophysics
%\item Jun-05	University of Sussex
%\item May-05	Universitaet Bielefeld, Germany
%\item Feb-05	Fermilab
%\item Jan-05	Oxford Astrophysics
%\item Nov-04	CERN Cosmology Seminar
%\item Oct-04	IFT Madrid
%\item May-04	CERN Cosmology Seminar
%\item Dec-03	Oxford Astrophysics
%\item Dec-03	Institute of Astronomy, Cambridge
%\item Oct-03	IAP
%\item Oct-03	Integral Data Center, Geneva
%\item Sep-03	University of Geneva
%\item May-03	Universita' La Sapienza, Roma
%\item May-03	Osservatorio Astronomico di Roma
%\item Mar-03	Univeristy of Zurich
%\item Jan-03	University of Geneva
%\item Oct-02	DAMPT, Cambridge
%\item May-02	University of Geneva
%\item Nov-01	IAS Princeton
%\item Jun-01	University of Geneva
%\end{enumerate}
\labelcvcount{nContribTalks}
\clearcvcounter
